\section{Related Work}
\label{sec:related}

\subsection{Anomaly Detection in Industrial Inspection}

Anomaly detection in visual data has been extensively studied, with methods broadly categorized into three paradigms: reconstruction-based, memory-bank-based, and vision-language model (VLM) approaches.

\textbf{Reconstruction-based methods} learn to reproduce only normal samples by training autoencoders~\cite{Bergmann2019MVTec}, variational autoencoders, or generative adversarial networks on defect-free images.
Anomalies are identified through reconstruction errors, under the assumption that the model cannot accurately reconstruct unseen defect patterns.
Recent extensions incorporate diffusion models~\cite{He2023FastRecon}, feature regression with distribution regularization, and perceptual losses to mitigate the common over-generalization problem where anomalous regions are inadvertently reconstructed well.

\textbf{Memory-bank-based methods} store representative patch features extracted from normal images and detect anomalies via nearest-neighbor retrieval in feature space.
SPADE~\cite{Cohen2022SPADE} pioneered this approach using multi-resolution feature correspondences.
PatchCore~\cite{Roth2022PatchCore} improved efficiency through coreset selection to reduce memory footprint while maintaining coverage.
AnomalyDINO~\cite{Deng2024AnomalyDINO} leverages DINOv2 features and demonstrates state-of-the-art few-shot performance through a sophisticated nearest-neighbor mechanism.
However, these methods typically require storing and searching through thousands to millions of patch descriptors, incurring substantial memory and computational overhead.

\textbf{VLM-based approaches} leverage vision-language models like CLIP~\cite{Radford2021CLIP} to enable text-guided anomaly detection.
WinCLIP~\cite{Huang2022WinCLIP} introduced manual text prompts with multi-scale windows for zero-shot and few-shot detection.
PromptAD~\cite{Li2024PromptAD} and IIPAD~\cite{Zhang2024IIPAD} automate prompt engineering or learn prompt adapters from normal instances.
While flexible, these methods often require auxiliary datasets, careful prompt design, or domain-specific textual priors.

\subsection{SubspaceAD: A Minimalist Baseline}

SubspaceAD~\cite{Lendering2025SubspaceAD} recently demonstrated that a remarkably simple pipeline can match or exceed these complex approaches.
It consists of only two stages: (1) extracting patch-level features from a small set of normal images using a frozen DINOv2-G backbone with multi-layer feature aggregation (layers 22--28), and (2) fitting a PCA model to estimate the low-dimensional linear subspace of normal variation.
At inference, anomaly scores are derived from the reconstruction residual with respect to this subspace.
Despite its simplicity, SubspaceAD achieves state-of-the-art results, with 1-shot image-level AUROC of 98.0\% on MVTec-AD and 93.3\% on VisA.

The success of SubspaceAD motivates our work: if a linear PCA subspace can achieve such strong performance, can a more expressive nonlinear model further improve anomaly discrimination?
Our MSAE-AD framework directly addresses this question by replacing PCA with a hierarchical sparse autoencoder that learns multi-granularity representations of normality.

\subsection{Sparse Autoencoders and MSAE}

Sparse autoencoders (SAEs) have emerged as powerful tools for discovering interpretable features in deep neural network representations~\cite{Bricken2023SAE,Templeton2024SAE}.
Traditional SAE variants---ReLU-based with L1 sparsity regularization and TopK-based with hard sparsity constraints---face a fundamental trade-off between reconstruction fidelity and sparsity.
ReLU SAEs achieve good reconstruction but struggle with sparsity control, while TopK SAEs enforce strict sparsity but suffer from dead neurons and degraded reconstruction quality.

Matryoshka SAE (MSAE)~\cite{Zaigrajew2025MSAE} was recently proposed to overcome this limitation.
MSAE introduces a nested (Matryoshka-style) architecture that learns hierarchical representations at multiple granularities simultaneously.
Instead of optimizing a single sparsity level, MSAE trains the autoencoder to reconstruct inputs from multiple subsets of latent neurons ordered by magnitude, enabling flexible sparsity control during inference through a weighting coefficient $\alpha$.
Two weighting strategies are proposed: uniform weighting (UW), where all granularity levels contribute equally, and reverse weighting (RW), which emphasizes higher-sparsity (fewer active neurons) reconstructions.
MSAE establishes a new state-of-the-art Pareto frontier between reconstruction quality and sparsity for interpreting CLIP embeddings, achieving 0.99 cosine similarity with 80\% sparsity.

Our work adapts MSAE from the interpretability domain to the anomaly detection task.
We hypothesize that MSAE's hierarchical sparse representations can capture the fine-grained structure of normal feature manifolds more effectively than a single linear subspace, leading to improved anomaly localization.
To the best of our knowledge, this is the first work to apply hierarchical sparse autoencoders to visual anomaly detection.
